\documentclass[12pt, letterpaper]{article}

%-------Packages---------
\usepackage{amssymb,amsfonts}
\usepackage{mathptmx}
\usepackage{amsthm}
\usepackage{graphicx}
\usepackage[all,arc]{xy}
\usepackage{enumerate}
\usepackage{bbm}
\usepackage{dsfont}
\usepackage{mathrsfs}
\usepackage{tikz-cd}
\usepackage{setspace}
\usepackage{import}
\usepackage[utf8]{inputenc}
\usepackage{hyperref}
\usepackage{tikz}
\usepackage{float}
\usepackage[dvipsnames]{xcolor}
\usepackage[nocheck]{fancyhdr}
\usepackage[nointegrals]{wasysym}

\usepackage[left=1in,top=1in,right=1in,bottom=1in]{geometry}

\usepackage{mathtools}
\usepackage{setspace}
\usepackage{cleveref}
\usepackage{multicol}
\usepackage{mathpazo}
\usepackage{tcolorbox}
\tcbuselibrary{theorems,skins,breakable}
\usepackage{xparse}

\usepackage{xcolor, ifthen, tcolorbox}
\tcbuselibrary{theorems,skins,breakable}
% Theorem System
% The following environments are provided:
%   New Problem:        \begin{problem} ...
%   New Question Header:   \qh (then followed by subproblems)
%   Subproblem:     \begin{subproblem} ...
%   Problem Proof:  \begin{problemproof} ...
%   Proof:          \\begin{pf}{color} ...
%   Remark:         \rmk (sentence), \rmkb (block)

% Define custom colors
\definecolor{thmscol}{HTML}{F4565B} %For Problems
\definecolor{lemscol}{HTML}{FE844C} %For Lemmes
\definecolor{prpscol}{HTML}{00A7EF} %For proposition
\definecolor{corscol}{HTML}{23DAFF} %For corrolaries
\definecolor{rmkscol}{HTML}{6DEEC5} %For remarks
\definecolor{defscol}{HTML}{18C991} %For definitions
\definecolor{exmscol}{HTML}{7DB800} %For examples
\definecolor{clmscol}{HTML}{165c58} %For claims

%Change between dark(black)/light(white) mode
\newcommand{\colormode}{white}
\ifthenelse{\equal{\colormode}{white}}
      {\newcommand{\TextColor}{black}}
      {\newcommand{\TextColor}{white}}
\newcommand{\pbcolormain}{thmscol!80!\colormode}
\newcommand{\pbcolorsecond}{thmscol!38!\colormode}


% ============================
% Problem Counter
% ============================
\newcounter{subproblemcounter}
\newcounter{problemcounter}

\newcommand{\q}{
    \setcounter{subproblemcounter}{0}%
    \stepcounter{problemcounter} % Increment problem counter
    \color{\TextColor}Problem \arabic{problemcounter}: % Display the problem counter
}

\newcommand{\stepsub}{
    \stepcounter{subproblemcounter}%
    \color{\TextColor}(\alph{subproblemcounter})
}
% ============================
% Problem
% ============================
\newtcolorbox{pb}[1]{
  enhanced,
  frame hidden,
  titlerule=0mm,
  toptitle=1mm,
  bottomtitle=1mm,
  fonttitle=\bfseries\large,
  coltitle=black,
  colbacktitle=\pbcolormain,
  colback=\pbcolorsecond,
  title={\q #1},
  coltext=\TextColor
}

\newtcolorbox{spb}[1]{
  enhanced,
  frame hidden,
  titlerule=0mm,
  toptitle=1mm,
  bottomtitle=1mm,
  fonttitle=\bfseries\large,
  coltitle=black,
  colbacktitle=\pbcolormain,
  colback=\pbcolorsecond,
  title={\stepsub #1},
  coltext=\TextColor,
  attach boxed title to top left={xshift=3mm,yshift=-2mm},
  boxed title style={
    colback=\pbcolormain,
    colframe=\pbcolormain,
  }
}

\newtcolorbox{pbh}[1]{
    enhanced,
    frame hidden,
    titlerule=0mm,
    toptitle=1mm,
    bottomtitle=1mm,
    fonttitle=\bfseries\large,
    coltitle=black,
    colbacktitle=\pbcolormain,
    colback=\pbcolorsecond,
    title={\q #1},
    boxrule=0pt,
    interior hidden,
    attach boxed title to top left={xshift=3mm,yshift=-2mm},
    boxed title style={
        colback=\pbcolormain,
        colframe=\pbcolormain,
    }
}

\newenvironment{problem}[1]{%
  \newpage
  \begin{pb}{#1}
    }{
  \end{pb}
}

\newenvironment{subproblem}[1]{%
  \begin{spb}{#1}
    }{
  \end{spb}
}

\newenvironment{problemproof}{%
    \begin{pf}{\pbcolormain}
        }{
    \end{pf}
}

\newcommand{\qh}[1][]{
    \newpage
    \begin{pbh}{#1}\end{pbh} % Display the problem counter
}

% ============================
% Claim
% ============================

\newtcbtheorem[number within=problemcounter]{claim}{Claim}
{
    enhanced,
    frame hidden,
    titlerule=0mm,
    toptitle=1mm,
    bottomtitle=1mm,
    fonttitle=\bfseries\large,
    coltitle=black,
    colbacktitle=clmscol!50!\colormode,
    colback=clmscol!20!\colormode,
}{clm}

\NewDocumentCommand{\clm}{m+m}{
    \begin{claim*}{#1}{}
        #2
    \end{claim*}
}

\newenvironment{clmpf}{
	{\noindent{\it \textbf{Proof.}}
	\setlength{\parindent}{0pt}
	}
	\tcolorbox[blanker,breakable,left=5mm,parbox=false,
    before upper={\parindent15pt},
    after skip=10pt,
	borderline west={1mm}{0pt}{clmscol!50!\colormode}]
}{
    \textcolor{clmscol!50!\colormode}{\hbox{}\nobreak\hfill$\blacksquare$}
    \endtcolorbox
}

\NewDocumentCommand{\clmp}{m+m+m}{
    \begin{claim*}{#1}{}
        #2
    \end{claim*}

    \begin{clmpf}
        #3
    \end{clmpf}
}


% ============================
% Proof
% ============================

\newenvironment{pf}[1]{%
  \def\pfcolor{#1}% Store the color in a macro
  \noindent{\it \textbf{Proof.}}%
  \tcolorbox[blanker,breakable,left=5mm,parbox=false,%
    before upper={\parindent15pt},%
    after skip=10pt,%
    borderline west={1mm}{0pt}{\pfcolor},%
    coltext=\TextColor
  ]%
  \setlength{\parindent}{0pt}%
}{%
  \textcolor{\pfcolor}{\hbox{}\nobreak\hfill$\blacksquare$}%
  \endtcolorbox%
}

\newtcolorbox{solution}{
  blanker,
  breakable,
  left=5mm,
  parbox=false,%
  before upper={\parindent15pt},%
  after skip=10pt,%
  borderline west={1mm}{0pt}{\pbcolormain},%
  coltext=\TextColor
}


% ============================
% Remark
% ============================


\NewDocumentCommand{\rmk}{+m}{
    {\it \color{rmkscol!100!white}#1}
}

\newenvironment{remark}{
    \par
    \vspace{5pt}
    \begin{minipage}{\textwidth}
        {\par\noindent{\textbf{Remark.}}}
        \tcolorbox[blanker,breakable,left=5mm,
        before skip=10pt,after skip=10pt,
        borderline west={1mm}{0pt}{rmkscol!80!\colormode},
        coltext=\TextColor
        ]
}{
        \endtcolorbox
    \end{minipage}
    \vspace{5pt}
}

\NewDocumentCommand{\rmkb}{+m}{
    \begin{remark}
        #1
    \end{remark}
}


% Define a new command for problems
\renewcommand{\headrule}{\color{\TextColor}\hrulefill}
\newcommand{\C}{\mathbb{C}}
\newcommand{\R}{\mathbb{R}}
\newcommand{\Q}{\mathbb{Q}}
\newcommand{\Z}{\mathbb{Z}}
\newcommand{\N}{\mathbb{N}}
\renewcommand{\P}{\mathbb{P}}
\newcommand{\F}{\mathbb{F}}
\newcommand{\T}{\mathcal{T}}
\newcommand{\contr}{\Rightarrow \Leftarrow}
\newcommand{\s}{\(\,\)}
\renewcommand{\t}[1]{\text{#1}}
\newcommand{\ang}[1]{\left\langle #1 \right\rangle}
\newcommand{\coker}{\mathrm{coker}}

% Meta data
\setstretch{1.2}

\begin{document}
\color{\TextColor}
\pagecolor{\colormode}
\setlength{\parindent}{0pt}
\pagestyle{fancy}
\fancyhf{}
\fancyhead[LOE]{\color{\TextColor}\slshape{\textbf{Vincent Ferrara}}}
\fancyhead[ROE]{\color{\TextColor}\thepage}

\newcommand{\mc}[1]{\mathcal{#1}}
\begin{problem}{}
    Let $X$ and $Y$ be topological spaces, and let $\mc{C}(X, Y)$ have
  the compact-open topology. Prove that if $Y$ is Hausdorff, then
  $\mc{C}(X, Y)$ is Hausdorff, and if $Y$ is regular, then
  $\mc{C}(X, Y)$ is regular.
\end{problem}

\clmp{}{
    \setlength{\parindent}{0pt}
    Let $X,Y$ be topological spaces, and let $\mathcal{C}(X,Y)$ have the compact-open topology. Let $U,V$ be open in $Y$.

    Prove that
  if $C$ is a compact set of $X$ and $\overline{U} \subset V$, then the
  subbasis sets satisfy $\overline{S(C, U)} \subset S(C, V)$.
}{
    \setlength{\parindent}{0pt}
    Assume $C$ is compact and $\overline{U} \subset V$. Let $f \in \overline{S(C,U)}$, and $x \in C$.

    Assume for contradiction that $f(x) \notin \overline{U}$. Then consider $W=Y \setminus \overline{U}$ which is open neighborhood of $f(x)$.

    Thus, consider $S(\{x\},W)$. Since $\{x\}$ is finite it is compact thus $S(\{x\},W)$ is an open neighborhood of $f$ since $f(x) \in W$.

    Thus, there exists a $g \in S(\{x\},W)$ s.t. $g \in S(C,U)$. However, this is a contradiction since this implies that $g(x) \in W \cap U \subseteq W \cap \overline{U} = \emptyset$.

    Therefore, $f(x) \in \overline{U} \subseteq V$. Thus, $f \in S(C,V)$.
}
\begin{problemproof}
    Assume $Y$ is Hausdorff. Let $f,g \in \mc{C}(X,Y)$ s.t. $f \neq g$. Thus, there is an $x \in X$ s.t. $f(x) \neq g(x)$. Then, because $Y$ is Hausdorff, you can
    find disjoint open neighborhoods $U$ and $V$ in $Y$ for $f(x)$ and $g(x)$.

    Now consider $\{x\}$ it is finite thus compact. Therefore, $S(\{x\},U)$ is an open neighborhood of $f$ and $S(\{x\},V)$ is an open neighborhood of $g$.

    Assume for contradiction that there is an $h \in S(\{x\},U) \cap S(\{x\},V)$. This implies that $h(x) \in U \cap V$ which is a contradiction. Thus, we have found disjoint neighborhoods of $f,g$. Thus, $\mathcal{C}(X,Y)$ is Hausdorff.

    Assume $Y$ is regular. Let $f \in \mathcal{C}(X,Y)$ and let $U$ be an open neighborhood of $f$. Thus, there exists subbasis sets s.t.
    \[f \in S(C_1,U_1) \cap S(C_2,U_2) \cap \dots \cap S(C_n,U_n) \subseteq U\]

    Let $y \in f(C_i)$, since $Y$ is regular this implies there is an open neighborhood $W_y$ s.t. $y \in \overline{W_y} \subset U_i$.

    Since $f(C_i)$ is compact and covered by these $W_y$'s. This implies there is a finite subcover s.t. $f(C_i) \subset W_{y_1} \cup \dots \cup W_{y_k}$.

    Define $W_i=W_{y_1} \cup \dots \cup W_{y_k}$.

    Since for all $j$ $\overline{W_{y_j}} \subset U_i$. This implies that $\overline{W_i} \subset U_i$.

    Now consider $S(C_i,W_i)$. Since $f(C_i) \subset W_i$ this implies that $f \in S(C_i,W_i)$.

    Therefore,

    \[f \in V=\bigcap_{i =1}^n S(C_i,W_i)\]

    Now consider $\overline{V}$

    \[\overline{V}=\overline{\bigcap_{i =1}^n S(C_i,W_i)} \subset \bigcap_{i =1}^n \overline{S(C_i,W_i)} \subset \bigcap_{i =1}^nS(C_i,U_i) \subset U\]

    Therefore, but the claim above $f \in \overline{V} \subset U$. Thus, $\mathcal{C}(X,Y)$ is regular.
\end{problemproof}

\begin{problem}{}
    Let $A_n$ be a sequence of matrices in $\textsc{GL}(d, \R)$, and let $A$ be
  another matrix. Recall that each $A_n$ and $A$ can be viewed as
  linear maps $\R^d \to \R^d$, and consider each of the four following
  conditions:
  \begin{enumerate}[(i)]
  \item The maps $A_n:\R^d \to \R^d$ converge \emph{pointwise} to $A$.
  \item The maps $A_n:\R^d \to \R^d$ converge \emph{uniformly} to $A$.
  \item The maps $A_n:\R^d \to \R^d$ converge \emph{uniformly on
      compacts} to $A$.
  \item For every $1 \le i,j \le d$, the $(i,j)$ entry of the matrix
    $A_n$ converges to the $(i,j)$ entry of $A$.
  \end{enumerate}

  Three of these conditions are equivalent to each other, but the
  fourth is not equivalent to any of the other three. Which one is not
  the same?
\end{problem}
\begin{problemproof}
    $(i),(iii),(iv)$ are equivalent.

    $(i) \implies (iv)$

    Assume $(i)$. Let $e_1,\dots,e_d$ be the standard basis vectors. This implies that

    \[A_n(e_j) \to A(e_j)\]

    Consider $\pi_i : \R^d \to \R$. Since this is continuous in $\R^d$ this implies that 
    
    $\pi_i(A_n(e_j))=(A_n)_{ij} \to A_{ij}$.

    Therefore, $(iv)$.

    $(iv) \to (iii)$
    
    Assume $(iv)$

    This implies that $|(A_n)_{ij}-A_{ij}| \to 0$.

    Now consider the norm

    \[|| A_n -A || = \sqrt{\sum_{i,j}|(A_n)_{ij}-A_{ij}|^2}\]

    Since there are finitely many times for any $\epsilon$ we just pick the max $N$ of all the terms thus $||A_n -A|| \to 0$.

    Now let $K$ be a compact set in $\R^n$. This means it is bounded. Thus, there is an $M > 0$ s.t. $||x|| < M$.

    Therefore,
    \[||(A_n-A)(x)|| \leq ||A_n-A||||x||\]

    So, for $\sup K$ this implies that 

    \[\sup_{x \in K} || (A_n-A)(x)|| \leq||A_n-A|| \sum_{x \in K}||x|| \leq M ||A_n-A||\]

    Since $||A_n-A|| \to 0 \implies \sup_{x \in K} || (A_n-A)(x)|| \to 0$.

    Therefore, $A_n$ converges uniformly on compacts to $A$.

    $(iii) \to (i)$

    Assume $(iii)$ compact convergence implies point wise convergence since $\{x\}$ is compact.

    Why $(ii)$ is not equivalent.

    Consider $A_n=(1+1/n)I_d$. $A_n$ converges point wise to the identity matrix. However, consider $\sup_{x \in \R^d}||A_n(x) -I_d(x)||$. Since $\R^d$ is unbounded consider a vector $x$ s.t. $||x||=M$. Thus, $||A_n(x)-I_d(x)||=(1/n)M$, since we can pick an $M \in \R$. This implies that the supremum is infinite and does not exist. Thus, it does not uniformly converge.
\end{problemproof}

\begin{problem}{}
    Let $X$ be a compact Hausdorff space, and equip the set $\t{Homeo}(X)$
  with the compact-open topology. Prove that $\t{Homeo}(X)$ is a
  topological group.
\end{problem}
\begin{problemproof}
    Let $G=\t{Homeo}(X)$
    
    Define $i : G \to G$ s.t. $i(x)=x^{-1}$.

    Let $C$ be compact and $U$ open in $X$. Consider $W=S(C,U) \cap G$ which is open in $G$.

    Consider $i^{-1}(W)$. 

    Consider $S(X \setminus U, X \setminus C) \cap G$. This is open since $X \setminus U$ is closed which implies compact, since $X$ is compact Hausdorff, and $X \setminus C$ is open because $C$ is compact in a Hausdorff space.

    Let $f \in i^{-1}(W)$ then $f^{-1}(C) \subset U \implies X \setminus U \subset f^{-1}(X \setminus C) \implies f(X \setminus U) \subset X \setminus C \implies f \in S(X \setminus U, X \setminus C)$.

    Let $f \in S(X \setminus U, X \setminus C) \implies f(X \setminus U) \subset X \setminus C \implies X \setminus U \subset f(X \setminus C) \implies f^{-1}(C) \subset U \implies f^{-1} \in W \implies f \in i^{-1}(W)$.

    Thus, $i^{-1}(W)=S(X \setminus U, X \setminus C) \cap G$ which is open in $G$. Thus, $i$ is continuous since we can build every other open set in $G$ from this type.

    Define $h : G \times G \to G$ s.t. $h(x,y)=x \circ y$.

    Consider $h^{-1}(W)$. Let $(f,g) \in h^{-1}(W)$. Thus, $f \circ g \in W$.

    Consider $g(C)$ which is compact and $f^{-1}(U)$ which is open.

    Since $g(C) \subset f^{-1}(U)$. Since $X$ is Hausdorff and compact thus locally compact. This implies that for $x \in g(C)$ there exists an open neighborhood $W_x$ of $x$ s.t. $\overline{W_x}$ is compact, and $\overline{W_x} \subset f^{-1}(U)$.

    Thus, $W_x$ form an open cover of $g(C)$ which is compact. Thus, there is a finite subcover $A=W_{x_1} \cup \dots \cup W_{x_n} \subset f^{-1}(U)$.

    Since $\overline{W_{x_i}} \subset f^{-1}(U) \implies \overline{A} \subset f^{-1}(U)$.

    Thus, consider $S(\overline{A},U) \times S(C, A)$.

    Let $(a,b) \in S(\overline{A},U) \times S(C, A)$. Then we can see that $a \circ b(C) \subset a(A) \subset a(\overline{A}) \subset U$.

    Thus, $a\circ b \in W$. Thus, $h$ is continuous.
\end{problemproof}

\begin{problem}{}
    Again, let $X$ be a compact Hausdorff space, and let
  $\mathrm{MCG}(X)$ denote the collection of connected components of
  $\t{Homeo}(X)$. One can define an operation $\cdot$ on
  $\mathrm{MCG}(X)$ as follows: if $C_1, C_2$ are elements in
  $\mathrm{MCG}(X)$, and $f_1 \in C_1$, $f_2 \in C_2$, then
  $C_1 \cdot C_2$ is the connected component of $\t{Homeo}(X)$ containing
  $f_1 \circ f_2$.

    Prove that the operation $\cdot$ makes $\mathrm{MCG}(X)$ into
    a group. \textit{Note:} most of the work here is proving that the
    operation $\cdot$ is actually well-defined, i.e. proving that it
    does not matter which maps $f_1, f_2$ you pick in the definition
    above.
\end{problem}
\begin{problemproof}
    Well-defined:

    Consider $C_1$ and $C_2$ which are connected components of Homeo$(X)$. Since these are both connected this implies that $C_1 \times C_2$ is connected in $\t{Homeo}(X) \times \t{Homeo}(X)$.

    Because the group operation map is continuous this implies that $h(C_1 \times C_2)$ is connected in $\t{Homeo}(X)$, so this is some subset of connected component lets call $C_3$.
    
    Thus, for $f,f' \in C_1$ and $g,g' \in C_2$, consider $h(f,g)=f\circ g \in C_3$ and $h(f',g')=f' \circ g' \in C_3$. Thus, the group operation is well defined.

    Associativity:

    This comes from function composition being associative.

    Identity:

    This is the connected component where the identity is because $[f] \cdot [id] = [f] = [id] \cdot [f]$.

    Inverses:

    For $f \in C_1$ the inverse is just which connected component $f^{-1}$ is in. Thus, $[f] \cdot [f^{-1}] = [id] = [f^{-1}] \cdot [f]$. 
\end{problemproof}

\begin{problem}{}
  Let $X$ be a completely regular space, and let $\beta(X)$ denote the
  Stone-\v{C}ech compactification of $X$. Prove that $X$ is connected
  if and only if $\beta(X)$ is connected.
\end{problem}
\begin{problemproof}
  Assume that $\beta(X)$ is disconnected.

  Thus, $\beta(X)=A \cup B$ for disjoint open sets $A,B$. Thus, $A \cap X$ and $B \cap X$ are disjoint open sets in $X$. These are non-empty since $X$ is dense in $\beta(X)$. Therefore,
  $X=A \cup B \cup X$. Thus, this implies that $X$ is disconnected.

  Assume that $X$ is disconnected. Let $X=A \cup B$ for disjoint open sets of $A,B$.

  Consider $f: X \to \{0,1\}$ s.t. $f(x)=1$ if $x \in A$ and 0 otherwise.

  Let $U$ be open in $\{0,1\}$. So $U=\{1\},\{0\},\{1,0\}$. In the first case $f^{-1}(\{1\})=A$, $f^{-1}(\{0\})=B$, $f^{-1}(\{0,1\})=X$. Thus, $f$ is continuous.

  Since $\{0,1\}$ is compact and Hausdorff this implies that $f$ can be extended to a continuous function $\hat{f} : \beta(X) \to \{0,1\}$ s.t. $\hat{f} \mid_X = f$.

  Also, $\hat{f}$ is still surjective since $\hat{f}\mid_X=f$ which is surjective. Thus, $\hat{f}^{-1}(0)$, and $\hat{f}^{-1}(1)$ are disjoint open sets that union to $\beta(X)$. Thus, $\beta(X)$ is disconnected.
\end{problemproof}
\end{document}